% fig_lec07_chart_basis.tex — Coordinate basis ∂_u Φ, ∂_v Φ along a curve
% Data: generated by scripts/sim_lec07_lec08.jl
%
% The coordinate basis vectors at points along a curve C(t) on M:
% they are tangent vectors to M, but they *change* from point to point.
% This motivates the need for Christoffel symbols — the chart-α partial
% derivative ∂_a tracks change of components, but the basis itself moves.

\begin{center}
\begin{tikzpicture}
  \begin{axis}[
      width=10.0cm, height=6.2cm,
      view={-32}{34},
      xmin=-3, xmax=3, ymin=-3, ymax=3, zmin=-1.1, zmax=1.1,
      axis lines=none,
      xtick=\empty, ytick=\empty, ztick=\empty,
      clip=false,
  ]
    % Manifold (faded mesh)
    \addplot3[surf, opacity=0.30, shader=interp,
              colormap name=grsurface, mesh/rows=55, forget plot]
      table {data/lec07_surface.dat};

    % Curve C(t)
    \addplot3[munsell, line width=1.3pt, no markers, forget plot]
      table {data/lec07_basis_curve.dat};

    % Sample points along the curve
    \addplot3[only marks, mark=*, mark size=1.8pt, banana!90!black, forget plot]
      table {data/lec07_basis_points.dat};

    % ∂_u Φ at each sample point
    \addplot3[
      quiver={u=\thisrowno{3}, v=\thisrowno{4}, w=\thisrowno{5}},
      -{Stealth[length=5pt, width=3.5pt]},
      banana!70!black, line width=1.0pt, no markers, forget plot,
    ] table[x index=0, y index=1, z index=2]
      {data/lec07_basis_du.dat};

    % ∂_v Φ at each sample point
    \addplot3[
      quiver={u=\thisrowno{3}, v=\thisrowno{4}, w=\thisrowno{5}},
      -{Stealth[length=5pt, width=3.5pt]},
      cgblue!75!black, line width=1.0pt, no markers, forget plot,
    ] table[x index=0, y index=1, z index=2]
      {data/lec07_basis_dv.dat};

    % Labels
    \node[font=\sf\small, text=spacecadet]
      at (axis cs:-2.4,-2.8,1.0) {$\M$};
    \node[font=\scriptsize, text=munsell, anchor=west]
      at (axis cs:1.55,0.45,0.30) {$C$};
    \node[font=\scriptsize, text=banana!55!black, anchor=south west]
      at (axis cs:-0.95,-0.25,0.55) {$\partial_u\Phi$};
    \node[font=\scriptsize, text=cgblue!70!black, anchor=south west]
      at (axis cs:-1.20,0.55,0.55) {$\partial_v\Phi$};
  \end{axis}
\end{tikzpicture}
\end{center}
